% Einheit 3 von 5 – Brüche multiplizieren und dividieren (bank/bruchrechnung/e3.jsonl, zusammenbau v0.1)
%% TODO Zweigzeile Teil 1: was hier gelernt wird (Satz in Schülersprache aus dem Einheitstitel) – Einheit 3
\einheitenkopf[e3]{Einheit 3 von 5 · Brüche multiplizieren und dividieren}
\zweigzeile{ab Kl. 6 (am Gymnasium ab Kl. 5) · keine P10-Aufgabe}
\setcounter{aufgabe}{17}

%% TODO Titel: Ich-kann-Satz fehlt in der Bank (Platzhalter: Kettenname) – Nr. 18
%% TODO Anweisung: ein Satz über der ersten Teilaufgabe fehlt in der Bank – Nr. 18
\begin{aufgabe}{Von heißt mal}
\begin{teile}
\teil Unterstreiche „von“ und schreibe die Malaufgabe auf: Zwei Fünftel von 35 Kindern fahren mit dem Rad. \leerfeld $\cdot$ \leerfeld
\end{teile}
\end{aufgabe}

%% TODO \verfahren{…}: Verfahrensüberschrift in Sachsprache fehlt in der Bank (2.3 b)
%% TODO Titel: Ich-kann-Satz fehlt in der Bank (Platzhalter: Kettenname) – Nr. 19
%% TODO Anweisung: ein Satz über der ersten Teilaufgabe fehlt in der Bank – Nr. 19
\begin{aufgabe}{Multiplizieren}
\begin{teile}
\teil „$\frac{1}{4}$ von 28 Kindern“ – heißt „von“ hier mal? \janein
\end{teile}
%% TODO Darstellung für schwach fehlt in der Bank (bruchrechnung-e3-k2-s1-v1; Abschnitt „Für schwache Schüler“)
\swz{$\frac{2}{9} \cdot 4$}{}
%% TODO Darstellung für schwach fehlt in der Bank (bruchrechnung-e3-k2-s1-v2; Abschnitt „Für schwache Schüler“)
\swz{$\frac{3}{10} \cdot 3$}{}
%% TODO Darstellung für schwach fehlt in der Bank (bruchrechnung-e3-k2-s1-v3; Abschnitt „Für schwache Schüler“)
\swz{$5 \cdot \frac{2}{11}$}{}
%% TODO Darstellung für schwach fehlt in der Bank (bruchrechnung-e3-k2-s1-v4; Abschnitt „Für schwache Schüler“)
\swz{$\frac{1}{8} \cdot 7$}{}
%% TODO Darstellung für schwach fehlt in der Bank (bruchrechnung-e3-k2-s1-v5; Abschnitt „Für schwache Schüler“)
\swz{$\frac{4}{15} \cdot 2$}{}
\swfrage{Was bleibt gleich, was ändert sich?}
%% TODO Darstellung für schwach fehlt in der Bank (bruchrechnung-e3-k2-s2-v1; Abschnitt „Für schwache Schüler“)
\swz{$\frac{2}{5} \cdot \frac{3}{7}$}{}
%% TODO Darstellung für schwach fehlt in der Bank (bruchrechnung-e3-k2-s3-v1; Abschnitt „Für schwache Schüler“)
\swz{$\frac{3}{8} \cdot \frac{4}{9}$}{}
%% TODO Darstellung für schwach fehlt in der Bank (bruchrechnung-e3-k2-s4-v1; Abschnitt „Für schwache Schüler“)
\swz{$1\frac{1}{2} \cdot \frac{4}{5}$}{}
\swz[3]{In einer Dose liegen 15 Murmeln: 7 grün, 5 weiß, 3 schwarz. Drei werden nacheinander ohne Zurücklegen gezogen. Tom behauptet: „Dreimal grün ist doppelt so wahrscheinlich wie dreimal weiß.“ Widerlege die Behauptung mit einer Rechnung. \hfill (P10 2019 OS)}{}
\swz[3]{Sechs Karten liegen verdeckt, eine davon ist die Gewinnkarte. Anna zieht zuerst eine Karte, danach Ben eine der übrigen fünf. Anna zieht keine Gewinnkarte mit der Wahrscheinlichkeit $\frac{5}{6}$, danach zieht Ben sie mit $\frac{1}{5}$. Wie groß ist die Wahrscheinlichkeit, dass Ben die Gewinnkarte zieht? \hfill (P10 2024 OS)}{}
\end{aufgabe}

%% TODO \verfahren{…}: Verfahrensüberschrift in Sachsprache fehlt in der Bank (2.3 b)
%% TODO Titel: Ich-kann-Satz fehlt in der Bank (Platzhalter: Kettenname) – Nr. 20
%% TODO Anweisung: ein Satz über der ersten Teilaufgabe fehlt in der Bank – Nr. 20
\begin{aufgabe}{Dividieren}
%% TODO Darstellung für schwach fehlt in der Bank (bruchrechnung-e3-k3-s1-v1; Abschnitt „Für schwache Schüler“)
\swz{$\frac{2}{3} : 5$}{}
%% TODO Darstellung für schwach fehlt in der Bank (bruchrechnung-e3-k3-s1-v2; Abschnitt „Für schwache Schüler“)
\swz{$\frac{5}{7} : 3$}{}
%% TODO Darstellung für schwach fehlt in der Bank (bruchrechnung-e3-k3-s1-v3; Abschnitt „Für schwache Schüler“)
\swz{$\frac{3}{8} : 4$}{}
%% TODO Darstellung für schwach fehlt in der Bank (bruchrechnung-e3-k3-s1-v4; Abschnitt „Für schwache Schüler“)
\swz{$\frac{7}{9} : 2$}{}
%% TODO Darstellung für schwach fehlt in der Bank (bruchrechnung-e3-k3-s1-v5; Abschnitt „Für schwache Schüler“)
\swz{$\frac{4}{5} : 3$}{}
\swfrage{Was bleibt gleich, was ändert sich?}
%% TODO Darstellung für schwach fehlt in der Bank (bruchrechnung-e3-k3-s2-v1; Abschnitt „Für schwache Schüler“)
\swz{$6 : \frac{1}{4}$ – wie oft passt $\frac{1}{4}$ in 6?}{}
%% TODO Darstellung für schwach fehlt in der Bank (bruchrechnung-e3-k3-s3-v1; Abschnitt „Für schwache Schüler“)
\swz{$\frac{1}{6} : \frac{3}{5}$}{}
%% TODO Darstellung für schwach fehlt in der Bank (bruchrechnung-e3-k3-s4-v1; Abschnitt „Für schwache Schüler“)
\swz{$\frac{5}{8} : \frac{15}{16}$}{}
\swz[3]{Ein Kanister enthält 9 Liter Apfelsaft. Er wird in Flaschen zu je $\frac{3}{4}$ Liter abgefüllt. Wie viele Flaschen werden voll?}{}
\end{aufgabe}

%% TODO Titel: Ich-kann-Satz fehlt in der Bank (Platzhalter: Kettenname) – Nr. 21
%% TODO Anweisung: ein Satz über der ersten Teilaufgabe fehlt in der Bank – Nr. 21
\begin{aufgabe}{Multiplizieren – Fehler finden, Anwendung, Darstellungswechsel}
\swz[3]{Emil rechnet: $\frac{2}{3} \cdot \frac{4}{5} = \frac{2}{3} \cdot \frac{5}{4} = \frac{10}{12}$. Finde den Fehler.}{}
\swz[3]{Ein Rezept für 4 Personen braucht $\frac{3}{4}$ Liter Milch. Du kochst für 6 Personen. Wie viel Milch brauchst du?}{}
\swa{Der Streifen hat 12 Teile. Färbe $\frac{1}{2}$ von $\frac{2}{3}$ des Streifens und gib den Anteil an. \leerfeld}{\bruchrechteck{0}{12}}
\end{aufgabe}

%% TODO Titel: Ich-kann-Satz fehlt in der Bank (Platzhalter: Kettenname) – Nr. 22
%% TODO Anweisung: ein Satz über der ersten Teilaufgabe fehlt in der Bank – Nr. 22
\begin{aufgabe}{Dividieren – Fehler finden, Begründen}
\swz[3]{Lars rechnet: $\frac{5}{8} : 4 = \frac{5}{2}$. Finde den Fehler.}{}
\swfrage{Warum ist $9 : \frac{1}{2}$ dasselbe wie $9 \cdot 2$?}
\end{aufgabe}

\uebersichtskasten{Mal: Zähler mal Zähler, Nenner mal Nenner – vorher kürzen, wenn es geht. „von“ heißt mal. \\ 2/3 · 5/8 = 10/24 = 5/12 \quad 3/4 von 20 = 3/4 · 20 = 15 \\ Geteilt: mit dem Kehrbruch malnehmen (Zähler und Nenner tauschen). \\ 5/6 : 2/3 = 5/6 · 3/2 = 15/12 = 1 1/4 \quad 4 : 1/2 = 4 · 2 = 8}
